\section{Toolbox: four definitions and the registry}\label{sec:toolbox}

This section holds all the theory the rest of the text needs: four definitions, then the registry, which contains the five lemmas that drive the outputs, three more arches, and the decks.

\subsection{Shapes and presheaves}

We call a small category $C$ a \emph{shape}. A \emph{presheaf} on $C$ is a functor $C^{\op}\to\Set$, in Lean \texttt{Cᵒᵖ ⥤ Type u}: a set $F(U)$ for each object, and for each arrow $f : V\to U$ a \emph{restriction} map $F(U)\to F(V)$, $x\mapsto\res xf$, with $\res x{\id} = x$ and $\res{(\res xg)}f = \res x{g\circ f}$. Presheaves on $C$ form the topos $\PSh(C)$. We use the following examples.

\begin{center}\small
\begin{tabularx}{\linewidth}{>{\raggedright\arraybackslash}p{4.3cm}>{\raggedright\arraybackslash}X}
\toprule
shape $C$ & presheaves on $C$\\
\midrule
$\B M$: one object, arrows $=$ elements of a monoid $M$, $g\circ f = gf$ & right $M$-sets: a set $X$ with $x\cdot m$, $x\cdot 1 = x$, $(x\cdot g)\cdot f = x\cdot gf$; in Mathlib, \texttt{MulAction Mᵐᵒᵖ X}\\
$\B(\Sigma^*)$, $\Sigma^*$ the words on an alphabet & deterministic automata: states $X$, transitions $x\cdot a$ (no initial or final states)\\
$\B(\mathbb N,+)$ & discrete dynamical systems: a set with a map $\tau : X\to X$\\
the free category on a quiver & representations of the opposite quiver: a set per vertex, a map per arrow\\
$V\rightrightarrows E$ & directed graphs\\
$\Delta$, $\Delta_+$ & simplicial sets, augmented simplicial sets\\
$\omega = (\mathbb N,\le)$ & sequences $X_0\leftarrow X_1\leftarrow\cdots$ (\q{the topos of trees} of guarded recursion)\\
\bottomrule
\end{tabularx}
\end{center}

For a monoid, the restriction of $x$ along $m$ is $x\cdot m$.

\subsection{Traps and negation}

A \emph{sub-presheaf} $A\subseteq F$, which we also call a \emph{trap}, is a family of subsets closed under restriction. For an $M$-set it is a subset closed under the action. Sub-presheaves are the subobjects in the topos and play the role of propositions. Their negation is not the set complement (which is usually not a trap) but the largest sub-presheaf disjoint from $A$:

\begin{lemma}\label{lem:neg}
$x\in\neg A$ iff no restriction of $x$ lies in $A$. Hence $x\in\neg\neg A$ iff every restriction of $x$ has a further restriction in $A$.
\end{lemma}

\begin{proof}
The set of such $x$ is closed under restriction, disjoint from $A$, and contains every trap disjoint from $A$. The second sentence is the first applied to $\neg A$.
\end{proof}

\begin{definition}
$\PSh(C)$ is \emph{Boolean} if $A\cup\neg A = F$ for every $A\subseteq F$ (excluded middle), and \emph{De Morgan} if $\neg A\cup\neg\neg A = F$ for every $A\subseteq F$ (weak excluded middle). By Lemma~\ref{lem:neg}, De Morgan means: \emph{for every presheaf $F$, every $A\subseteq F$ and every element $x$, either no restriction of $x$ lies in $A$, or every restriction of $x$ has a further restriction in $A$.}
\end{definition}

Informally, De Morgan is a \emph{stability} property: \q{$A$ can be reached} cannot be lost by restricting further. Boolean is stronger: \q{$x\in A$} is decided.

\subsection{The Ore condition}

$C$ satisfies the \emph{Ore condition} if every cospan $f : Y\to X$, $g : Z\to X$ completes to a commutative square $f\circ h = g\circ k$. For $\B M$: for all $a,b\in M$ there are $c,d$ with $ac = bd$, i.e.\ $aM\cap bM\neq\emptyset$ (any two elements have a common right multiple). It holds for groups, commutative monoids, monoids with a zero, categories with an initial object or with pullbacks, and chains.

\subsection{The registry}\label{sec:registry}

\subsubsection{The five arches}

The library \lp{} proves the following five bridges in Lean~4. We quote the main declarations, and give the exact statements of the first two.

\begin{bridge}[Boolean]\label{br:bool}
$\PSh(C)$ is Boolean iff $C$ is a groupoid.
\begin{leancode}
theorem isBooleanTopos_presheaf_iff_isGroupoid (C : Type u) [SmallCategory C] :
    IsBooleanTopos (Cᵒᵖ ⥤ Type u) ↔ IsGroupoid C
\end{leancode}
\end{bridge}

\begin{bridge}[De Morgan]\label{br:dm}
$\PSh(C)$ is De Morgan iff $C$ satisfies the Ore condition.
\begin{leancode}
theorem isDeMorganTopos_presheaf_iff_rightOreCondition (C : Type u)
    [SmallCategory C] :
    IsDeMorganTopos (Cᵒᵖ ⥤ Type u) ↔ GrothendieckTopology.RightOreCondition C
\end{leancode}
Its proof is short: if $\res xf\in A$ and $g$ is another arrow into $X$, complete the square $f\circ h = g\circ k$; then $\res{(\res xg)}k = \res{(\res xf)}h\in A$. Conversely, test on the presheaf $\y X = \Hom(-,X)$, the trap generated by $f$, and the element $\id_X$.
\end{bridge}

\begin{bridge}[Morita]\label{br:morita}
$\PSh(C)\simeq\PSh(D)$ iff the \emph{Karoubi envelopes} $\Kar(C)$ and $\Kar(D)$ are equivalent. $\Kar(C)$ adds an object for each idempotent arrow $e = e\circ e$; for $\B M$ its objects are the idempotents of $M$, and $e\cong f$ in $\Kar(\B M)$ iff there are $a\in fMe$, $b\in eMf$ with $ba = e$ and $ab = f$. Lean: \lean{moritaEquivalentCat_iff_karoubi}; corollaries \lean{isGroupoid_iff_of_moritaEquivalentCat}, \lean{rightOreCondition_iff_of_moritaEquivalentCat}.
\end{bridge}

\begin{bridge}[Essential points]\label{br:pts}
The functors $\PSh(C)\to\Set$ that preserve all limits and all colimits (the \emph{essential points}, or \q{fibre functors}) form a category equivalent to $\Kar(C)^{\op}$. For $\B M$: they are the functors $X\mapsto Xe = \{x : x\cdot e = x\}$ for $e$ idempotent, two of them being isomorphic iff $e\cong f$ in $\Kar(\B M)$. Lean: \lean{essentialPointEquivKaroubiOp}, \lean{isEssentialPoint_iff}.
\end{bridge}

\begin{bridge}[Subtopoi]\label{br:sub}
The subtopoi of $\PSh(C)$ (full subcategories closed under limits whose inclusion has a finite-limit-preserving left adjoint) correspond, order-reversingly, to the \emph{Grothendieck topologies} on $C$: choices of \q{covering} sieves on each object, stable under pullback and transitive. Lean: \lean{subtoposOrderIso}, \lean{exists_orderIso_ici_of_sheafEquiv}.
\end{bridge}

Every statement below is one of these five lemmas, specialised to a shape and unfolded with Lemma~\ref{lem:neg}.

\subsubsection{Three more arches}\label{sec:morearches}

The registry has three more arches. So far we have applied them only to the decks taken from Grothendieck's notes, below, and to his pseudo-fields (\S\ref{sec:pseudofields}); none of the outputs of Sections~\ref{sec:machines} to~\ref{sec:shapes} comes from them.

\begin{bridge}[Two-valued]\label{br:twoval}
The subterminal objects of $\PSh(C)$ (its \q{opens}) correspond to the sieves of $C$ \cite[IV~8.4.4]{sga4}. Hence $\PSh(C)$ has no subterminal objects besides $0$ and $1$ iff for any two objects $a$, $b$ of $C$ there is an arrow $a\to b$. For $C$ empty the topos is degenerate and the condition holds vacuously, so the statement is about nonempty $C$. Lean: \lean{forall_subfunctor_terminal_eq_bot_or_eq_top_iff}.
\end{bridge}

\begin{bridge}[Connected]\label{br:conn}
$\PSh(C)$ is connected and nonempty in the sense of \cite[IV~8.7~e)]{sga4} (every partition of its final object has exactly one nonempty part) iff $C$ is connected, a fact recalled as known in \cite[Ex.~3.6]{hora-kamio}. In terms of subterminal objects: $C$ is nonempty and the only complemented subterminal objects are $0$ and $1$. Lean: \lean{isConnected_iff_forall_isCompl}.
\end{bridge}

\begin{bridge}[Spatial]\label{br:spatial}
$\PSh(C)$ is equivalent to the presheaf topos of a preorder, and so to the topos of sheaves on its Alexandrov space, iff $C$ is thin. Grothendieck states it as an exercise \cite[p.~26]{grothendieck-161-3}; see \S\ref{sec:precedents}. Lean: \lean{exists_preorder_moritaEquivalentCat_iff_isThin}, which states the preorder form; that the topos is then the topos of a space is Grothendieck's equivalence, from the same exercise, and is not formalised here.
\end{bridge}

\subsubsection{The decks}

The deck registry holds the shapes of the table above, with their presheaves, and the shapes introduced with each output in \S\ref{sec:machines}--\ref{sec:shapes}. Step~2 of \S\ref{sec:method} explains how to recognise a new deck.

\subsubsection{The decks of Grothendieck's notes}\label{sec:grdecks}

Three decks recur in Grothendieck's working notes. Here is what the three arches above say about them (Lean: \lean{GrothendieckDecks.lean}).
\begin{itemize}
\item The simplex category $\Delta$, the first test category of \emph{Pursuing Stacks} (folders 108 and 118 of the fonds): simplicial sets form a two-valued and connected topos, since there is a constant map between any two objects of $\Delta$; this topos is not the presheaf topos of a preorder, since the two vertices of $[1]$ are two maps $[0]\to[1]$ (\lean{simplexCategory_twoValued}, \lean{simplexCategory_isConnected}, \lean{simplexCategory_not_spatial}).
\item A group $G$ as a one-object groupoid, whose presheaf topos is $\B G$ (folder 161-6, p.~22, where Grothendieck reduces a question to this case): the topos is two-valued (as \cite[IV~8.4.3]{sga4} notes for any monoid) and connected, and it is the presheaf topos of a preorder iff $G$ is trivial (\lean{singleObj_twoValued}, \lean{singleObj_isConnected}, \lean{singleObj_spatial_iff}).
\item A preordered set $I$, whose presheaf topos $\widehat I$ is $\mathrm{Top}(I)$ (folder 37, p.~48; folder 161-3, p.~26): the topos is two-valued iff any two elements are comparable both ways, that is, iff its posetal quotient has at most one point (\lean{preorder_twoValued_iff}).
\end{itemize}
These are immediate corollaries. We record them because they are the first entries of the registry taken from the fonds.
